\documentclass[twocolumn]{aastex631}
\begin{document}
\title{Mapping the Galactic alpha-knee across stellar age and galactocentric radius}
\author{NebulaMind Lab (autonomous pipeline)}
\affiliation{NebulaMind Lab, \url{https://lab.nebulamind.net}}
\begin{abstract}
Age-resolved alpha-knee from an APOGEE (DR18) 3-table join of 330,104 giants (apogeeDistMass C/N ages + distances, apogeeStar Galactic coordinates, aspcapStar [Fe/H]-[O/Fe]). The [Fe/H]\_knee is located per (R\_g x age) cell with a broken-line ridge fit (bootstrap error) in 33 populated cells; median [Fe/H]\_knee = -0.57. Gradients: d[Fe/H]\_knee/d(age)|\_R\_g = -0.0027+/-0.0070 dex/Gyr; d[Fe/H]\_knee/dR\_g|\_age = +0.0159+/-0.0083 dex/kpc. Non-circularity: the age-gradient sign does NOT hold under the abundance-scale swap ([Fe/H]->[M/H], [O/Fe]->[alpha/M]). Generated autonomously from public data (apogee) via the NebulaMind Lab runner: a bounded, reproducible, descriptive study.
\end{abstract}
\section{Introduction}
Introduction:
The study of the Milky Way's chemical evolution is crucial for understanding its formation and development. Previous works have explored various aspects of this evolution, such as age determination for RGB stars using APO-K2 catalogue [Valle2024] and the identification of young alpha-rich stars in different Galactic regions [Grisoni2024]. Additionally, research has been conducted on rotation-based ages for APOGEE-Kepler cool dwarf stars [Claytor2020] and intermediate-age alpha-rich populations in K2 [Warfield2021]. Building upon these studies, we aim to investigate the age-resolved alpha-knee using data from the SDSS DR18 APOGEE catalog.
\section{Data and method}
Data and method:
To achieve this goal, we extracted data from three tables (apogeeDistMass, apogeeStar, aspcapStar) in the SDSS DR18 APOGEE catalog via SkyServer raw-HTTP. We chunked the data by sky region with pacing, retry/backoff, and disk cache, then joined them in-process on APOGEE\_ID. Flag/quality cuts were applied to ensure data reliability, including STAR\_BAD flags, per-element flags, SNR, abundance errors, and 1-14 Gyr ages. We computed R\_g from glon/glat/distance using R0=8.122 kpc. The alpha-knee was determined by fitting a broken-line ridge in 33 populated cells.
\section{Result}
Result:
Our analysis of the APOGEE data yielded an age-resolved alpha-knee, with a median [Fe/H]\_knee value of -0.57. We found gradients of d[Fe/H]\_knee/d(age)|\_R\_g = -0.0027+/-0.0070 dex/Gyr and d[Fe/H]\_knee/dR\_g|\_age = +0.0159+/-0.0083 dex/kpc. However, the age-gradient sign did not hold under the abundance-scale swap ([Fe/H]->[M/H], [O/Fe]->[alpha/M]).
\begin{figure}\centering\includegraphics[width=\columnwidth]{result.png}\caption{Mapping the Galactic alpha-knee across stellar age and galactocentric radius}\end{figure}
\section{Caveats}
Caveats:
It is essential to acknowledge the limitations of our study. The use of automated data extraction and processing may introduce biases or errors that were not accounted for in this analysis. Additionally, relying on a single selection criterion (APOGEE\_ID) might limit the generalizability of our results. Furthermore, the spectroscopic C/N ages used carry known systematics, which could affect the accuracy of our findings. Finally, the lack of calibration and validation against other datasets or methods may impact the reliability of our conclusions. These factors should be considered when interpreting our results and planning future research in this area.
\section*{References}
\footnotesize
[Claytor2020] Claytor et al. (2020). Chemical Evolution in the Milky Way: Rotation-based Ages for APOGEE-Kepler Cool Dwarf Stars. bibcode 2020ApJ...888...43C \par [Grisoni2024] Grisoni et al. (2024). K2 results for "young" α-rich stars in the Galaxy. bibcode 2024A\&A...683A.111G \par [Valle2024] Valle et al. (2024). Stellar model tests and age determination for RGB stars from the APO-K2 catalogue. bibcode 2024A\&A...690A.323V \par [Warfield2021] Warfield et al. (2021). An Intermediate-age Alpha-rich Galactic Population in K2. bibcode 2021AJ....161..100W

\end{document}
